\section*{الفصل \ref{ch:g8:the-atmosphere} --- الغلاف الجوي}

\begin{solution}{exo:g8:the-atmosphere:1}
(أ) \omterm{def:g8:the-atmosphere:layers}{الستراتوسفير}؛ (ب) \omterm{def:g8:the-atmosphere:layers}{التروبوسفير}؛ (ج) \omterm{def:g8:the-atmosphere:layers}{الميزوسفير}؛ (د) \omterm{def:g8:the-atmosphere:layers}{الثرموسفير}.
\end{solution}

\begin{solution}{exo:g8:the-atmosphere:2}
(أ) \qty{15}{\degreeCelsius}. (ب) نحو \qty{-56}{\degreeCelsius}.
(ج) نحو \qty{-3}{\degreeCelsius}. (د) نحو \qty{-86}{\degreeCelsius}.
(ه) من \qty{11}{km} إلى نحو \qty{20}{km} (وعلى مدى بضعة كيلومترات
قرب \qty{50}{km} وفوق \qty{86}{km} مباشرة).
\end{solution}

\begin{solution}{exo:g8:the-atmosphere:3}
(أ) خطأ: فهي تنخفض في \omterm{def:g8:the-atmosphere:layers}{التروبوسفير}، وترتفع في \omterm{def:g8:the-atmosphere:layers}{الستراتوسفير}، وتنخفض في
\omterm{def:g8:the-atmosphere:layers}{الميزوسفير} وترتفع من جديد في \omterm{def:g8:the-atmosphere:layers}{الثرموسفير}. (ب) صواب.
(ج) خطأ: إنها في \omterm{def:g8:the-atmosphere:layers}{الستراتوسفير}. (د) صواب.
\end{solution}

\begin{solution}{exo:g8:the-atmosphere:4}
(أ) \omterm{def:g8:the-atmosphere:layers}{التروبوسفير}. (ب) \omterm{def:g8:the-atmosphere:layers}{الستراتوسفير} (ففي الغلاف الجوي المعياري يقع
\omterm{def:g8:the-atmosphere:layers}{التروبوبوز} على \omterm{def:g4:world-physical-map:altitude}{ارتفاع} \qty{11}{km}). (ج) \omterm{def:g8:the-atmosphere:layers}{الستراتوسفير}. (د) \omterm{def:g8:the-atmosphere:layers}{الميزوسفير}.
\end{solution}

\begin{solution}{exo:g8:the-atmosphere:5}
(أ) $20 - 6.5 \times 2 = 20 - 13 = \qty{7}{\degreeCelsius}$. (ب) $20 -
6.5 \times 3 = 20 - 19.5 = \qty{0.5}{\degreeCelsius}$. (ج) يجب أن يبرد
الهواء بمقدار 20 درجة: $20 \div 6.5 \approx 3.1$، أي على \omterm{def:g4:world-physical-map:altitude}{ارتفاع} نحو
\qty{3.1}{km}.
\end{solution}

\begin{solution}{exo:g8:the-atmosphere:6}
(أ) $3 - 0.5 = \qty{2.5}{km}$. (ب) $6.5 \times 2.5 = 16.25$: نحو 16
درجة. (ج) $18 - 16.25 \approx \qty{2}{\degreeCelsius}$.
\end{solution}

\begin{solution}{exo:g8:the-atmosphere:7}
(أ) \qty{500}{hPa} و\qty{250}{hPa} و\qty{125}{hPa}. (ب) قرب
الأرض القاعدة جيدة (\qty{500}{hPa} مقابل \qty{506}{hPa}). أما في الأعلى
فتعطي أكثر مما ينبغي: \qty{250}{hPa} مقابل \qty{227}{hPa}،
و\qty{125}{hPa} مقابل \qty{96}{hPa}. فقاعدة التنصيف ليست إلا دليلًا
تقريبيًا.
\end{solution}

\begin{solution}{exo:g8:the-atmosphere:8}
(أ) \babelsublr{$10~\text{hPa} = 10^1$~hPa}؛ $10^3 \div 10^1 = 10^2$: أصغر مئة
مرة. (ب) \babelsublr{$0.01~\text{hPa} = 10^{-2}$~hPa}. (ج) $10^3 \div 10^{-2} =
10^5$: أصغر مئة ألف مرة.
\end{solution}

\begin{solution}{exo:g8:the-atmosphere:9}
(أ) $6 \div 1.5 = 4$ مرات أعرض. (ب) كلما ارتفع، ضغط عليه الهواء \omterm{def:g2:water-on-earth:ocean}{المحيط} به
أقل فأقل، فيتمدد الغاز الذي في داخله ويكبر البالون حتى ينشق غلافه. (ج)
بعد انفجار البالون، تبطئ المظلة سقوط المسبار اللاسلكي، كي لا يؤذي أحدًا
ولا يتلف شيئًا.
\end{solution}

\begin{solution}{exo:g8:the-atmosphere:10}
(أ) $20 \div 7 \approx 2.9$ مرة أسمك. (ب) أقرب إلى \omterm{def:g3:continents-and-oceans:poles}{القطبين}:
\omterm{def:g4:world-physical-map:altitude}{فارتفاع} \qty{11}{km} يزيد بمقدار \qty{4}{km} على \qty{7}{km} لكنه يقل بمقدار \qty{9}{km} عن \qty{20}{km} (فهو يصف العروض الوسطى، لا \omterm{def:g3:continents-and-oceans:equator}{خط الاستواء}).
(ج) في البداية يكون \omterm{def:g8:the-atmosphere:layers}{التروبوبوز} على \omterm{def:g4:world-physical-map:altitude}{ارتفاع} نحو \qty{7}{km}، تحت
الطائرة: فالطائرة في \omterm{def:g8:the-atmosphere:layers}{الستراتوسفير}. وفي النهاية يكون \omterm{def:g8:the-atmosphere:layers}{التروبوبوز} على \omterm{def:g4:world-physical-map:altitude}{ارتفاع}
نحو \qty{20}{km}: فالطائرة في \omterm{def:g8:the-atmosphere:layers}{التروبوسفير}.
\end{solution}

\begin{solution}{exo:g8:the-atmosphere:11}
(أ) $1 - 506 \div 1013 \approx 0.50$: \textbf{50~\%}. (ب) $1 - 227 \div
1013 \approx 0.78$: \textbf{78~\%}. (ج) $1 - 12 \div 1013 \approx 0.99$:
\textbf{99~\%}. (د) نحو 1~\%.
\end{solution}

\begin{solution}{exo:g8:the-atmosphere:12}
(أ) $6371 \div 15 \approx 425$~كيلومترًا لكل سنتيمتر. (ب) $100 \div 425
\approx 0.24$~سنتيمتر، أي نحو \qty{2.4}{mm}. (ج) $11 \div 425 \approx
0.026$~سنتيمتر، أي نحو \qty{0.26}{mm}. (د) $0.26 \div 0.1 \approx 2.6$: ورقتان
أو ثلاث ورقات.
\end{solution}

\begin{solution}{pb:g8:the-atmosphere:1}
\textbf{1.} $300 \times 10 = 3000$~متر.

\textbf{2.} $6000 \div 300 = 20$ دقيقة؛ $\num{30000} \div 300 = 100$
دقيقة.

\textbf{3.} \babelsublr{100} دقيقة، أي ساعة واحدة و40~دقيقة.

\textbf{4.} على \omterm{def:g4:world-physical-map:altitude}{ارتفاع} \qty{8}{km}: \omterm{def:g8:the-atmosphere:layers}{التروبوسفير} (فدرجة الحرارة لا تزال
تنخفض). وعلى \omterm{def:g4:world-physical-map:altitude}{ارتفاع} \qty{30}{km}: \omterm{def:g8:the-atmosphere:layers}{الستراتوسفير}.

\textbf{5.} $15 - (-50) = 65$ درجة.

\textbf{6.} $65 \div 10 = \qty{6.5}{\degreeCelsius}$ لكل كيلومتر.

\textbf{7.} من \qty{12}{km} إلى \qty{20}{km} تبقى عند
\qty{-56}{\degreeCelsius}. ويتوقف التبرد بين \qty{10}{km}
و\qty{12}{km}: \omterm{def:g8:the-atmosphere:layers}{فالتروبوبوز} على \omterm{def:g4:world-physical-map:altitude}{ارتفاع} نحو \qty{11}{km}.

\textbf{8.} ترتفع: \qty{-55}{\degreeCelsius} على \omterm{def:g4:world-physical-map:altitude}{ارتفاع} \qty{22}{km}،
و\qty{-47}{\degreeCelsius} على \omterm{def:g4:world-physical-map:altitude}{ارتفاع} \qty{30}{km}. فالبالون في
\omterm{def:g8:the-atmosphere:layers}{الستراتوسفير}، الذي يسخّنه الأوزون الذي يمتص الأشعة فوق البنفسجية الآتية
من الشمس.

\textbf{9.} $1013 \div 2 \approx \qty{507}{hPa}$؛ ويُبلغ بين
\qty{4}{km} (\qty{617}{hPa}) و\qty{6}{km} (\qty{472}{hPa}).

\textbf{10.} $1013 \div 4 \approx \qty{253}{hPa}$.

\textbf{11.} بين \qty{10}{km} (\qty{265}{hPa}) و\qty{12}{km}
(\qty{194}{hPa}).

\textbf{12.} من \qty{10}{km}، يجب أن ينخفض الضغط بمقدار \babelsublr{$265 - 253 =
12$~hPa} أخرى، من أصل \babelsublr{71~hPa} يفقدها على مدى كيلومترين: $12 \div 71 \times 2
\approx 0.34$~كيلومتر. فيبلغ الضغط ربع قيمته عند \omterm{def:g4:world-physical-map:altitude}{مستوى سطح البحر} على
\omterm{def:g4:world-physical-map:altitude}{ارتفاع} نحو \textbf{\qty{10.3}{km}}.
\end{solution}
